% Lecture notes: Week 8, Lecture 2 -- Beam Deflection: Superposition
% To hide this lecture from the website, add a line containing only:  % draft
% To set the date shown on the website, add a line like:  % posted: 2026-09-02
\documentclass[11pt]{article}
\usepackage{mech2030notes}

\begin{document}

% ##################################################################
%  WEEK 8, LECTURE 2
% ##################################################################
\lecture{8}{2}{Beam Deflection: Superposition}

\noindent
$\rightarrow$ Consider:

\noindent
\begin{minipage}[c]{0.6\textwidth}
\centering
\begin{tikzpicture}
  % scale: 0.6 cm per m
  \simplebeam{6}
  \distload{3}{6}{0.15}{0.7}{5}
  \node[above] at (4.5,0.85) {$\SI{4}{kN/m}$};
  \draw[force] (3,1.6) node[above] {$\SI{20}{kN}$} -- (3,0.17);
  \node[pin] at (3,0) {}; \node[below] at (3,-0.15) {$C$};
  \node[left] at (-0.4,0.1) {$A$};
  \node[right] at (6.4,0.1) {$B$};
  \draw[dim] (0,-1.3) -- (3,-1.3) node[midway, fill=white] {$\SI{5}{m}$};
  \draw[dim] (3,-1.3) -- (6,-1.3) node[midway, fill=white] {$\SI{5}{m}$};
\end{tikzpicture}
\end{minipage}%
\hfill
\begin{minipage}[c]{0.36\textwidth}
(Steel)
\[
  \left\{
  \begin{aligned}
    E &= \SI{200}{GPa} \\
    I &= \SI{0.1457e-3}{m^4}
  \end{aligned}
  \right.
\]
\end{minipage}

\bigskip
\noindent
\textbf{Q:} What is $v_C$?

\section*{Principle of Superposition}

\noindent
$\rightarrow$ The \emph{principle of superposition} says we can sum the contributions due to
the two loads separately.
\begin{itemize}[label=\then]
  \item True for \emph{small displacements} and \emph{linear elasticity}.
\end{itemize}

\subsection*{Equivalent:}

\begin{center}
\begin{tikzpicture}[baseline=(b)]
  \coordinate (b) at (0,0);
  \simplebeam{5}
  \draw[force] (2.5,1.4) node[above] {$\SI{20}{kN}$} -- (2.5,0.17);
  \node[pin] at (2.5,0) {};
  \draw[force] (2.5,-0.2) -- (2.5,-1.1) node[below] {$v_C^{1}$};
\end{tikzpicture}
\qquad{\Large $+$}\qquad
\begin{tikzpicture}[baseline=(b)]
  \coordinate (b) at (0,0);
  \simplebeam{5}
  \distload{2.5}{5}{0.15}{0.7}{4}
  \node[above] at (3.75,0.85) {$\SI{4}{kN/m}$};
  \node[pin] at (2.5,0) {};
  \draw[force] (2.5,-0.2) -- (2.5,-1.1) node[below] {$v_C^{2}$};
\end{tikzpicture}
\end{center}

\subsection*{1) From tables}

\begin{align*}
  v_C^{1} &= -\frac{PL^3}{48EI}
           = -\frac{(\SI{20e3}{N})(\SI{10}{m})^3}{48\,(\SI{200e9}{Pa})(\SI{0.1457e-3}{m^4})}
\end{align*}
\[
  \boxed{v_C^{1} = -\SI{14.3}{mm}}
\]

\subsection*{2) From tables}

\begin{align*}
  v_C^{2} &= -\frac{5wL^4}{768EI}
           = -\frac{5\,(\SI{4e3}{N/m})(\SI{10}{m})^4}{768\,(\SI{200e9}{Pa})(\SI{0.1457e-3}{m^4})}
\end{align*}
\[
  \boxed{v_C^{2} = -\SI{8.94}{mm}}
\]

\subsection*{Superpose the two}

\[
  v_C = v_C^{1} + v_C^{2}
  \qquad\Rightarrow\qquad
  \fbox{\fbox{$v_C = -\SI{23.24}{mm}$}}
\]

\noindent
$\rightarrow$ Superposition also applies to angles, e.g.
\begin{align*}
  \theta_A &= \theta_A^{1} + \theta_A^{2} \\
  \theta_B &= \theta_B^{1} + \theta_B^{2}
\end{align*}

\end{document}
